\worksheettitle
  {Разряд показывает значение}
  {\modulemark{M01-R} Коррекционная карточка}

\studentnameblock

\begin{notebox}[Как выполнять]
  Если вы называете саму цифру вместо её значения или считаете разряды
  слева, выполните карточку по шагам.
\end{notebox}

\section{Двигайтесь справа налево}

\begin{center}
\renewcommand{\arraystretch}{1.5}
\begin{tabular}{c|ccccc}
  \toprule
  разряд & дес. тысяч & ед. тысяч & сотни & десятки & единицы \\
  \midrule
  цифра в числе \(81\,804\) & 8 & 1 & 8 & 0 & 4 \\
  значение & \answerblank[18mm] & \answerblank[16mm] &
    \answerblank[14mm] & \answerblank[14mm] & \answerblank[14mm] \\
  \bottomrule
\end{tabular}
\end{center}

\begin{notebox}[Проговорите]
  Первая цифра \(8\) означает восемь десятков тысяч, то есть \(80\,000\).
  Вторая цифра \(8\) означает восемь сотен, то есть \(800\).
\end{notebox}

\section{Сравните}

В числе \(35\,353\):

\begin{enumerate}[label=\textbf{\arabic*.}]
  \item Значение первой цифры \(3\): \answerblank[22mm].
  \item Значение второй цифры \(3\): \answerblank[22mm].
  \item Значение третьей цифры \(3\): \answerblank[22mm].
\end{enumerate}

\begin{checkpointbox}[Повторная проверка]
  Рассмотрите новое число \(90\,609\).

  Назовите разряды двух цифр \(9\):
  \answerblank[32mm] и \answerblank[32mm].

  Назовите значения этих цифр:
  \answerblank[25mm] и \answerblank[25mm].

  Я объясняю разницу так:
  \writinglines{1}
\end{checkpointbox}
